\documentclass[a4paper]{article}

\usepackage{graphicx}
\usepackage{amsmath,amssymb}
\usepackage{minted}
\usepackage{multirow}
\usepackage{float}
\usepackage{todonotes}
\usepackage{forest}
\usepackage{xstring}
\usepackage{xcolor}
\usepackage[most]{tcolorbox}
\usepackage{xparse}
\usepackage{natbib}

\usetikzlibrary{graphs,graphdrawing}
\usegdlibrary{layered}

\input{/home/donkarlo/Dropbox/repo/nd_ai_project/src/nd_ai/knoweldge_representation/language/formal/computer/programming/tex/latex/code_clips/external_graphviz.tex}
\input{/home/donkarlo/Dropbox/repo/nd_ai_project/src/nd_ai/knoweldge_representation/language/formal/computer/programming/tex/latex/parts/control_sequence/kind/semtag/semtag.tex}
\input{/home/donkarlo/Dropbox/repo/nd_ai_project/src/nd_ai/knoweldge_representation/language/formal/computer/programming/tex/latex/code_clips/seehere.tex}

\title{}
\date{}

\begin{document}
    \maketitle
    
    
    \section{Präsens}
        \textbf{Verwendung}: Gegenwart, allgemeine Aussagen, Zukunft mit Zeitangabe.\\
        \textbf{Form}: Subjekt + konjugiertes Verb.\\
        ex: Ich lerne Deutsch: I learn German.\\
        ex: Er arbeitet heute: He works today.\\
    
    
    \section{Perfekt}
        \textbf{Verwendung}: gesprochene Vergangenheit.\\
        \textbf{Form}: Subjekt + haben/sein + Partizip II.\\
        ex: Ich habe Deutsch gelernt: I have learned German.\\
        ex: Er ist nach Hause gegangen: He has gone home.\\
    
    
    \section{Präteritum}
        \textbf{Verwendung}: schriftliche Vergangenheit.\\
        \textbf{Form}: Subjekt + Präteritum.\\
        ex: Ich lernte Deutsch: I learned German.\\
        ex: Er ging nach Hause: He went home.\\
    
    
    \section{Plusquamperfekt}
        \textbf{Verwendung}: Handlung vor einer Vergangenheit.\\
        \textbf{Form}: Subjekt + hatte/war + Partizip II.\\
        ex: Ich hatte gegessen, bevor er kam: I had eaten before he came.\\
        ex: Sie war gegangen: She had left.\\
    
    
    \section{Futur I}
        \textbf{Verwendung}: Zukunft, Vermutung.\\
        \textbf{Form}: Subjekt + werden + Infinitiv.\\
        ex: Ich werde lernen: I will learn.\\
        ex: Er wird zu Hause sein: He is probably at home.\\
    
    
    \section{Futur II}
        \textbf{Verwendung}: abgeschlossene Zukunft, Vermutung über Vergangenheit.\\
        \textbf{Form}: Subjekt + werden + Partizip II + haben/sein.\\
        ex: Ich werde gelernt haben: I will have learned.\\
        ex: Er wird angekommen sein: He will have arrived.\\
    
    
    \section{Konjunktiv I}
        
        \subsection{Indirekte Rede}
            \textbf{Verwendung}: indirekte Aussagen.\\
            \textbf{Form}: Subjekt + Konjunktiv I.\\
            ex: Er sagt, er sei krank: He says he is sick.\\
            ex: Sie sagt, sie habe Zeit: She says she has time.\\
    
    
    \section{Konjunktiv II}
        
        \subsection{Irrealer Wunsch}
            \textbf{Verwendung}: unrealer Wunsch.\\
            \textbf{Form}: würde + Infinitiv oder Präteritumform.\\
            ex: Ich würde reisen: I would travel.\\
            ex: Ich hätte mehr Zeit: I would have more time.\\
        
        \subsection{Bedingung}
            \textbf{Verwendung}: hypothetische Bedingung.\\
            \textbf{Form}: Wenn + Konjunktiv II, würde + Infinitiv.\\
            ex: Wenn ich Zeit hätte, würde ich kommen: If I had time, I would come.\\
            ex: Wenn er reich wäre, würde er ein Haus kaufen: If he were rich, he would buy a house.\\
        
        \subsection{Modalverben im Konjunktiv II}
            \textbf{Wichtig}: Modalverben haben eigene Formen (ohne würde).\\
            \textbf{können} $\rightarrow$ könnte\\
            \textbf{müssen} $\rightarrow$ müsste\\
            \textbf{dürfen} $\rightarrow$ dürfte\\
            \textbf{sollen} $\rightarrow$ sollte\\
            \textbf{wollen} $\rightarrow$ wollte\\
            \textbf{mögen} $\rightarrow$ möchte\\
            ex: Ich könnte kommen: I could come.\\
            ex: Ich müsste arbeiten: I would have to work.\\
            ex: Ich dürfte gehen: I would be allowed to go.\\
            ex: Ich sollte mehr lernen: I should learn more.\\
            ex: Ich wollte gehen: I wanted to go (oder höflich: would want).\\
            ex: Ich möchte einen Kaffee: I would like a coffee.\\
        
        \subsection{Höflichkeit}
            \textbf{Verwendung}: höfliche Fragen und Bitten.\\
            ex: Könnten Sie mir helfen: Could you help me.\\
            ex: Würden Sie das erklären: Would you explain that.\\
        
        \subsection{Irreale Vergangenheit}
            \textbf{Form}: hätte/wäre + Partizip II.\\
            ex: Ich hätte mehr gelernt: I would have learned more.\\
            ex: Er wäre gekommen: He would have come.\\
    
    
    \section{Imperativ}
        \textbf{Verwendung}: Befehle und Bitten.\\
        \textbf{Form}: Verb im Imperativ.\\
        ex: Komm her: Come here.\\
        ex: Gehen Sie bitte: Please go.\\
    
    
    \section{Passiv}
        
        \subsection{Präsens}
            \textbf{Form}: werden + Partizip II.\\
            ex: Das Auto wird repariert: The car is being repaired.\\
        
        \subsection{Präteritum}
            \textbf{Form}: wurde + Partizip II.\\
            ex: Das Auto wurde repariert: The car was repaired.\\
        
        \subsection{Perfekt}
            \textbf{Form}: ist + Partizip II + worden.\\
            ex: Das Auto ist repariert worden: The car has been repaired.\\
        
        \subsection{Plusquamperfekt}
            \textbf{Form}: war + Partizip II + worden.\\
            ex: Das Auto war repariert worden: The car had been repaired.\\
        
        \subsection{Futur I}
            \textbf{Form}: wird + Partizip II + werden.\\
            ex: Das Auto wird repariert werden: The car will be repaired.\\
    
    
    \section{Zusatz: Aktiv vs Passiv}
        \textbf{Aktiv}: Der Mann repariert das Auto.\\
        \textbf{Passiv}: Das Auto wird repariert.\\

\end{document}